\section{Related Work}
\label{sec:related}

\subsection*{Weight-Space Self-Supervised Learning}

The idea of learning representations from neural network weights ---
treating checkpoints as data --- has attracted growing attention.
Hyper-Representations~\cite{Schurholt2021Hyper} first demonstrated that training an autoencoder
on populations of neural network weights produces latent codes that recover hyperparameters,
accuracy, and generalization gap without task labels.
SANE~\cite{Schurholt2024Towards} extended this to a hierarchical transformer architecture
with contrastive SSL, achieving $R^2 = 0.72$ on heterogeneous model zoo property prediction
within its intended same-architecture setting.
Both methods use \emph{flat tokenization}: weight matrices are chunked into fixed-size tokens,
processed positionally, and aggregated ---
an approach that implicitly assumes weight token positions are comparable across networks.

This assumption holds when training and test architectures share the same layer structure.
When architectures differ substantially (e.g., MLP/CNN vs.\ ViT), flat tokenization loses meaning.
Our work identifies the consequence:
SANE's latent representations collapse to near-constant codes for cross-architecture inputs
($\text{std} \approx 0.0014$), and property prediction drops to $R^2 = 0.072$.
This is a cross-architecture distribution shift failure, not a failure of SANE as a method
in its intended domain.
The graph encoding approach we adopt directly addresses this structural mismatch.

\subsection*{Equivariant Graph Neural Networks for Weight Encoding}

A parallel research strand represents neural networks as computation graphs and uses
equivariant GNNs as encoders.
The foundational theoretical work~\cite{Navon2023Equivariant} characterized all
permutation-equivariant/invariant linear layers over weight spaces.
Neural Functional Transformers~\cite{Zhou2023Permutation} and
NF-Layers~\cite{Kofinas2024Graph} demonstrated that graph-based encoders achieve strong
performance on property prediction within a fixed architecture family.
Critically, Kofinas et al.~\cite{Kofinas2024Graph} showed that graph-based encoders can
generalize \emph{across} architecture families in the \emph{supervised} setting,
achieving $R^2 = 0.71$ on held-out CNNs after training on MLPs.
ScaleGMN~\cite{Kalogeropoulos2024Scale} extended this to include scale equivariance,
showing $+8$--$12$ $R^2$ points over permutation-only baselines in supervised same-architecture
settings.
Our work differs in two key dimensions:
(1) we operate in the \emph{self-supervised} setting; and
(2) we target \emph{cross-architecture} transfer to ViTs with LayerNorm,
which changes the benefit profile of scale equivariance.

\subsection*{Cross-Architecture Weight Representation Transfer}

The most directly related work is Ballerini et al.~\cite{Ballerini2025Weight},
who applied contrastive SSL to diverse neural radiance field architectures and demonstrated
cross-architecture property transfer without target-domain training data.
Our work generalizes this paradigm to general-purpose neural architectures
(MLPs, CNNs, ViTs) --- a broader and more architecturally diverse setting.
The ViT Model Zoo~\cite{Falk2025ModelZoo}, which provides real ViT-S/16 ImageNet
checkpoints with accuracy labels, enables evaluation of cross-architecture SSL at scale;
we use their dataset as our test set.

\subsection*{Symmetry and Normalization Layers}

The interaction between equivariance symmetry groups and neural network normalization has
received recent theoretical attention.
We hypothesize that LayerNorm in transformer architectures fixes the gauge degree of freedom
corresponding to activation scale, reducing the effective symmetry group of ViT weight spaces
to permutation-only.
Our empirical finding --- that scale equivariance \emph{hurts} ViT cross-architecture transfer ---
provides experimental support for this hypothesis.

Existing weight-space SSL methods achieve strong results within architecture families but
fail under architectural distribution shift.
Existing equivariant graph encoders achieve cross-architecture transfer in supervised settings
but do not address the SSL setting.
We bridge these, discovering that scale equivariance's benefit is
architecture-normalization-dependent.
